\subsubsection{Cubical geometric semantics of time fibers: hyperplane statistics, symmetry, and bulk limits}\label{subsubsec:pom-fiber-cubical-hyperplane-symmetry-limit}
In \S\ref{subsubsec:pom-fiber-reconstruction-cubical}, the fiber reconstruction graph \(\mathcal F_x\) has been established as a class of auditable cubical graph objects:
under the \((011\leftrightarrow 100)\) local kernel, its structure is entirely determined by the component lengths \((\ell_i)\) of the invertible-site induced graph \(\Gamma_x=\bigsqcup_i P_{\ell_i}\),
and satisfies
\(
\mathcal F_x\cong\prod_i \Gamma_{\ell_i}
\)
(Theorem \ref{thm:pom-fiber-reconstruction-cartesian-product}).
Therefore, the auditable objects of time fibers (preimage fibers) are not merely the cardinality \(d_m(x)=|\Fold_m^{-1}(x)|\), but also include the hyperplane (\(\Theta\)-class) statistics, symmetry, and limiting probability structures carried by the \(1\)-skeleton of the CAT(0) cubical complex.
This subsection presents a collection of results that are self-contained within this semantics; the inputs are solely the component lengths \((\ell_i)\) and the Fibonacci cube \(\Gamma_{\ell}\) corresponding to each component (Definition \ref{def:pom-fibonacci-cube}).

\paragraph{Notational convention (to avoid confusion with parameter slices)}
In this subsection, \(\Theta\) denotes the \(\Theta\)-equivalence class (hyperplane) induced by the Djokovi\'c--Winkler relation of a partial cube,
and is not to be conflated with the ``finite audit slice'' notation \(\Theta\) in \S\ref{subsec:pom-three-generator-residual-closure}; the two contexts are disjoint and never appear simultaneously in the same expression.

\paragraph{Symmetry: rigidity of the automorphism group}

\begin{theorem}[Automorphism group of the Fibonacci cube]\label{thm:pom-fibcube-aut-z2}
For all \(n\ge 2\), the graph automorphism group of the Fibonacci cube \(\Gamma_n\) satisfies
\[
\mathrm{Aut}(\Gamma_n)\ \cong\ \ZZ_2,
\]
whose nontrivial element may be taken as the graph automorphism induced by coordinate reversal (index reflection) \(b_1\cdots b_n\mapsto b_n\cdots b_1\).
\end{theorem}
\begin{proof}
For each edge of \(\Gamma_n\), the \(\Theta\)-relation partitions the edges into \(n\) \(\Theta\)-classes; denote by \(\Theta_i\) the set of edges formed by flipping the \(i\)-th coordinate (see Theorem \ref{thm:pom-fibcube-theta-class-size} below).
Two classes \(\Theta_i,\Theta_j\) \emph{cross} in the cubical complex sense if and only if there exists a \(4\)-cycle (a sub-cube \(Q_2\)) whose two pairs of opposite edges lie in \(\Theta_i\) and \(\Theta_j\), respectively.
In \(\Gamma_n\), if \(|i-j|\ge 2\), then at the all-zero word one can independently flip positions \(i\) and \(j\) to obtain such a \(Q_2\); if \(|i-j|=1\), any simultaneous flip forces an adjacent \(11\), making such a \(Q_2\) impossible.
Hence the ``non-crossing relation'' on \(\{\Theta_1,\dots,\Theta_n\}\) induces a path graph \(P_n\) (adjacent if and only if \(|i-j|=1\)).
\par
Any graph automorphism \(\varphi\in\mathrm{Aut}(\Gamma_n)\) must permute the set of \(\Theta\)-classes while preserving the ``crossing/non-crossing'' relation,
thereby inducing an automorphism of \(P_n\); this gives a group injection
\(
\mathrm{Aut}(\Gamma_n)\hookrightarrow \mathrm{Aut}(P_n)\cong \ZZ_2
\).
On the other hand, index reversal clearly preserves the ``no adjacent \(11\)'' constraint and preserves Hamming distance \(1\), thus providing a nontrivial automorphism.
Therefore \(\mathrm{Aut}(\Gamma_n)\cong\ZZ_2\). This completes the proof.
\end{proof}

\begin{theorem}[Automorphism group of the general fiber reconstruction graph]\label{thm:pom-fiber-reconstruction-aut-group}
If \(\mathcal F_x\cong\prod_{i=1}^r\Gamma_{\ell_i}\), and we write
\(
m_\ell:=\#\{\,i:\ \ell_i=\ell\,\}
\),
then the automorphism group of \(\mathcal F_x\) admits the canonical decomposition
\[
\mathrm{Aut}(\mathcal F_x)\ \cong\
\Bigl(\prod_{i=1}^r \ZZ_2\Bigr)\rtimes
\Bigl(\prod_{\ell}\mathfrak S_{m_\ell}\Bigr),
\]
so that
\(
|\mathrm{Aut}(\mathcal F_x)|=2^{r}\prod_{\ell}(m_\ell!)
\).
\end{theorem}
\begin{proof}
For Cartesian products of connected graphs, the prime factor decomposition is unique up to isomorphism, and the automorphism group is generated by ``per-factor automorphisms \(\times\) permutations of isomorphic factors'' (a standard result in graph product theory).
By Theorem \ref{thm:pom-fiber-reconstruction-cartesian-product}, the Cartesian decomposition of \(\mathcal F_x\) is canonically fixed by \((\ell_i)\);
by Theorem \ref{thm:pom-fibcube-aut-z2}, the automorphism group of each factor \(\Gamma_{\ell_i}\) is isomorphic to \(\ZZ_2\),
and Fibonacci cubes \(\Gamma_\ell\) of different lengths cannot be isomorphic since their vertex counts \(F_{\ell+2}\) differ.
Therefore, only factors of equal length may be permuted, yielding \(\prod_\ell\mathfrak S_{m_\ell}\);
the per-factor inversions yield \(\prod_i\ZZ_2\);
combining the two gives the stated semidirect product decomposition and order formula. This completes the proof.
\end{proof}

\begin{corollary}[The component length multiset is a complete isomorphism invariant]\label{cor:pom-fiber-reconstruction-spectrum-complete}
If \(\mathcal F_x\cong \mathcal F_y\) (graph isomorphism), then
\(
\{\ell_i(x)\}_i
\)
and
\(
\{\ell_j(y)\}_j
\)
are equal as multisets.
Therefore, the fiber spectrum
\[
\mathsf{Spec}_{\square}(x):=\{\ell_1,\dots,\ell_r\}
\]
serves as a complete isomorphism invariant of \(\mathcal F_x\); this spectrum simultaneously controls \(|\mathrm{Aut}(\mathcal F_x)|\) (Theorem \ref{thm:pom-fiber-reconstruction-aut-group}).
\end{corollary}
\begin{proof}
By Theorem \ref{thm:pom-fiber-reconstruction-cartesian-product}, the prime factor decompositions (in the Cartesian product sense) of \(\mathcal F_x\) and \(\mathcal F_y\) yield two collections of Fibonacci cube factors.
The uniqueness of prime factorization implies that the two factor multisets coincide;
since the vertex count \(F_{\ell+2}\) of \(\Gamma_\ell\) is strictly increasing in \(\ell\), factor isomorphism holds if and only if the lengths are equal, yielding multiset equality of \((\ell_i)\). This completes the proof.
\end{proof}

\paragraph{Hyperplane statistics: \(\Theta\)-class sizes and \(k\)-point functions}

\begin{theorem}[Closed form for \(\Theta\)-class edge counts]\label{thm:pom-fibcube-theta-class-size}
View \(\Gamma_n\subseteq\{0,1\}^n\) as a partial cube, and for each coordinate \(1\le i\le n\) denote by \(\Theta_i\) the edge set obtained by flipping the \(i\)-th coordinate (i.e., the \(i\)-th \(\Theta\)-class/hyperplane).
Let the Fibonacci sequence satisfy \(F_0=0,F_1=1\).
Then for any \(n\ge 1\) and \(1\le i\le n\),
\[
|\Theta_i|\ =\ F_i\,F_{n-i+1}.
\]
In particular,
\[
|\Theta_1|=|\Theta_n|=F_n,\qquad
|\Theta_i|=F_n-F_{i-1}F_{n-i}\ (<F_n)\quad(1<i<n),
\]
so the endpoint hyperplanes are uniquely characterized as the ``largest \(\Theta\)-classes.''
\end{theorem}
\begin{proof}
Each \(\Theta_i\)-edge connects exactly one pair of vertices differing at position \(i\); one endpoint has position \(i\) equal to \(1\).
That endpoint and the edge determine each other uniquely (flipping position \(i\) from \(1\) to \(0\) gives the other endpoint),
so \(|\Theta_i|\) equals the number of vertices in \(\Gamma_n\) with position \(i\) equal to \(1\).
For a valid word of length \(n\), the condition \(v_i=1\) forces \(v_{i-1}=v_{i+1}=0\) (boundaries handled by convention),
so the free portion splits into left and right segments: the left segment of length \(i-2\) has \(F_{(i-2)+2}=F_i\) valid words, and the right segment of length \(n-i-1\) has \(F_{(n-i-1)+2}=F_{n-i+1}\) valid words.
Their product gives \(|\Theta_i|=F_iF_{n-i+1}\).
\par
The second identity follows from the Fibonacci identity \(F_n=F_iF_{n-i+1}+F_{i-1}F_{n-i}\); when \(1<i<n\), the second term is positive, so the count is strictly less than \(F_n\).
This completes the proof.
\end{proof}

\begin{corollary}[Fibonacci ratio closed form and boundary layer exponential decay for coordinate marginals]\label{cor:pom-fibcube-marginal-boundary-layer}
Let \(\varphi=\frac{1+\sqrt5}{2}\).
Under the uniform counting measure on \(\Gamma_n\), for any \(1\le i\le n\) we have the strict closed form
\[
\mathbb P(v_i=1)=\frac{|\Theta_i|}{|V(\Gamma_n)|}
=\frac{F_iF_{n-i+1}}{F_{n+2}}.
\]
Moreover, when \(n\to\infty\) with \(i\to\infty\) and \(n-i\to\infty\), the marginal probability converges to the Parry bulk constant
\[
\mathbb P(v_i=1)\ \longrightarrow\ \pi(1)=\frac{1}{\varphi^2+1}=\frac{1}{\varphi\sqrt5}
\qquad\text{(see Theorem \ref{thm:pom-parry-limit-chain-explicit})},
\]
with exponentially decaying deviation
\[
\mathbb P(v_i=1)-\frac{1}{\varphi^2+1}
=O\!\left(\varphi^{-2i}+\varphi^{-2(n-i+1)}\right)
\qquad(n\to\infty).
\]
More generally, if \(\mathcal F_x\cong\prod_{j=1}^r\Gamma_{\ell_j}\), then for any coordinate \((j,a)\) we have
\[
\mathbb P(v_{j,a}=1)=\frac{F_aF_{\ell_j-a+1}}{F_{\ell_j+2}},
\]
with the same exponential boundary layer decay
\(
O(\varphi^{-2a}+\varphi^{-2(\ell_j-a+1)})
\)
in the bulk window \(a\to\infty\), \(\ell_j-a\to\infty\).
\end{corollary}
\begin{proof}
The first identity follows immediately from Theorem \ref{thm:pom-fibcube-theta-class-size} and \(|V(\Gamma_n)|=F_{n+2}\).
The limit and error bound follow by substituting the Binet expansion
\(
F_k=\frac{\varphi^k}{\sqrt5}\bigl(1+O(\varphi^{-2k})\bigr)
\)
into the closed form and simplifying.
The product fiber case reduces to single-factor estimates via the tensorization of Proposition \ref{prop:pom-fiber-hyperplane-tensorization}. \qed
\end{proof}

\begin{theorem}[Exact count of coordinates equal to \(1\) and the general \(k\)-point function]\label{thm:pom-fibcube-kpoint-count}
Let \(V(\Gamma_n)=X_n\) denote the set of valid words of length \(n\), equipped with uniform counting.
For any index set \(S=\{i_1<\cdots<i_k\}\subseteq\{1,\dots,n\}\), if there exists \(j\) such that \(i_{j+1}-i_j=1\), then
\[
\#\{v\in V(\Gamma_n):\ v_{i_j}=1\ \forall j\}=0.
\]
If \(i_{j+1}-i_j\ge 2\) holds (the necessary and sufficient condition for simultaneous occupancy), then there is a strict product decomposition
\[
\#\{v\in V(\Gamma_n):\ v_{i_j}=1\ \forall j\}
=
F_{i_1}\cdot
\Bigl(\prod_{j=1}^{k-1}F_{i_{j+1}-i_j-1}\Bigr)\cdot
F_{n-i_k+1}.
\]
\end{theorem}
\begin{proof}
When adjacent indices exist, the condition forces an adjacent \(11\), so the count is \(0\).
When the gap condition holds, fixing \(v_{i_j}=1\) simultaneously forces each neighboring position to \(0\), splitting the length-\(n\) valid word into \(k+1\) mutually disjoint free segments:
the left free segment has length \(i_1-2\), the \(j\)-th middle segment has free length \(i_{j+1}-i_j-3\), and the right free segment has length \(n-i_k-1\).
The number of valid words in each segment is \(F_{i_1}\), \(F_{i_{j+1}-i_j-1}\), and \(F_{n-i_k+1}\), respectively;
inter-segment independence yields the product. This completes the proof.
\end{proof}

\paragraph{Bulk limit: Parry measure, alternating correlations, and CLT}

\begin{theorem}[Fully explicit Parry limit chain]\label{thm:pom-parry-limit-chain-explicit}
Let \(X_\infty\subseteq\{0,1\}^{\mathbb N}\) be the golden-mean SFT (forbidding adjacent \(11\)), and take its Parry maximum entropy equilibrium state \(\nu_{\mathrm{Parry}}\).
Then \(\nu_{\mathrm{Parry}}\) is a two-state Markov chain with transition matrix
\[
P=\begin{pmatrix}
\varphi^{-1} & \varphi^{-2}\\
1 & 0
\end{pmatrix},
\qquad \varphi=\frac{1+\sqrt5}{2}.
\]
Its stationary distribution is
\[
\pi(1)=\frac{1}{\varphi^2+1},\qquad
\pi(0)=\frac{\varphi^2}{\varphi^2+1},
\]
and the unique nontrivial eigenvalue is \(\lambda_2=-\varphi^{-2}\).
\end{theorem}
\begin{proof}
This is the standard Parry construction given by the Perron--Frobenius normalization of the golden-mean adjacency matrix \(A=\bigl(\begin{smallmatrix}1&1\\1&0\end{smallmatrix}\bigr)\) (see also Remark \ref{prop:parry-measure-golden}).
Direct computation verifies that \(P\) is stochastic with stationary distribution \(\pi\); the eigenvalues follow from \(\det(P-\lambda I)=0\), which apart from \(\lambda_1=1\) yields only \(\lambda_2=-\varphi^{-2}\). This completes the proof.
\end{proof}

\begin{proposition}[Bulk limit: local statistics of uniform valid words converge to the Parry measure]\label{prop:pom-unif-xn-local-limit-parry}
For each \(n\ge 1\), let \(U_n:=\mathrm{Unif}(X_n)\) denote the uniform measure on the set of valid words of length \(n\).
Fix any window length \(k\ge 1\) and any valid block \(a_1\cdots a_k\in\{0,1\}^k\) (no adjacent \(11\)).
If the position \(t=t(n)\) satisfies \(t\to\infty\) and \(n-t-k\to\infty\), then the local limit holds:
\[
U_n\bigl(X_t=a_1,\dots,X_{t+k-1}=a_k\bigr)\ \longrightarrow\
\nu_{\mathrm{Parry}}\bigl([a_1\cdots a_k]\bigr),
\]
where the right-hand side is the cylinder set probability under the Parry measure (Theorem \ref{thm:pom-parry-limit-chain-explicit}).
\end{proposition}
\begin{proof}
Express the valid word count using the golden-mean adjacency matrix \(A=\bigl(\begin{smallmatrix}1&1\\1&0\end{smallmatrix}\bigr)\) and its Perron root \(\varphi\):
\(|X_n|=\mathbf 1^{\mathsf T}A^{n-1}\mathbf 1\) (boundary conventions do not affect the exponential scale).
For a fixed block \(a_1\cdots a_k\), its count at position \(t\) decomposes as ``left completion count \(\times\) intra-block transitions \(\times\) right completion count,''
which is a product of certain power matrix entries of \(A\) and constant factors.
As \(t\to\infty\) and \(n-t-k\to\infty\),
by the Perron--Frobenius theorem,
\(
A^m=\varphi^m(vw^{\mathsf T}+o(1))
\)
(\(v,w\) being the right and left Perron vectors) holds entry-wise;
substituting into the above decomposition and taking the ratio with \(|X_n|\), the \(o(1)\) error vanishes simultaneously on both sides,
yielding the limiting distribution as precisely the Parry measure cylinder set probability normalized by \(v,w\),
which is equivalent to the Markov expression in Theorem \ref{thm:pom-parry-limit-chain-explicit}. This completes the proof.
\end{proof}

\begin{corollary}[Two-point function and alternating correlation length]\label{cor:pom-parry-two-point-alternating}
Under \(\nu_{\mathrm{Parry}}\), let \(X_t\in\{0,1\}\) denote the value at position \(t\).
Then for any \(k\ge 0\), the exact formula holds:
\[
\mathbb P(X_0=1,X_k=1)
=\pi(1)^2+\pi(0)\pi(1)\,(-\varphi^{-2})^{k},
\]
so that the covariance satisfies
\[
\mathrm{Cov}(X_0,X_k)=\pi(0)\pi(1)\,(-\varphi^{-2})^{k},
\]
exhibiting strictly alternating sign and exponential decay rate \(\varphi^{-2}\).
\end{corollary}
\begin{proof}
By the Markov property, \(\mathbb P(X_0=1,X_k=1)=\pi(1)\,(P^k)_{11}\).
For the two-state chain, the spectral decomposition of \(P^k\) has only one component besides the stationary direction, namely \(\lambda_2^k\), and the coefficient can be determined from \(\mathbb P(X_0=1,X_1=1)=0\) at \(k=1\),
yielding the stated closed form; the covariance is obtained by subtracting \(\pi(1)^2\). This completes the proof.
\end{proof}

\begin{theorem}[CLT for the count of \(1\)'s in valid words: \(\varphi\)-closed-form mean and variance densities]\label{thm:pom-parry-ones-clt}
Under \(\nu_{\mathrm{Parry}}\), let
\(
S_n:=\sum_{t=1}^{n}X_t
\).
Then \(S_n\) satisfies the central limit theorem
\[
\frac{S_n-n\mu}{\sqrt n}\ \Longrightarrow\ \mathcal N(0,\sigma^2),
\qquad
\mu=\pi(1)=\frac{1}{\varphi^2+1},
\]
where the variance density admits the closed form
\[
\sigma^2
=\pi(0)\pi(1)\frac{1+\lambda_2}{1-\lambda_2}
=\frac{\varphi^3}{(\varphi^2+1)^3}.
\]
\end{theorem}
\begin{proof}
For ergodic finite-state reversible Markov chains, any bounded additive observable satisfies a CLT;
here \(X_t\) is a bounded function of the chain, so the conclusion holds.
The variance density is given by the Green--Kubo formula:
\(
\sigma^2=\mathrm{Var}(X_0)+2\sum_{k\ge 1}\mathrm{Cov}(X_0,X_k)
\).
Substituting the covariance from Corollary \ref{cor:pom-parry-two-point-alternating} and summing yields
\[
\sigma^2=\pi(0)\pi(1)\left(1+2\sum_{k\ge 1}\lambda_2^k\right)
=\pi(0)\pi(1)\frac{1+\lambda_2}{1-\lambda_2},
\]
and substituting \(\pi\) and \(\lambda_2=-\varphi^{-2}\) gives the closed form. This completes the proof.
\end{proof}

\begin{remark}[General finite-window additive observables]\label{rem:pom-parry-additive-functional-clt}
The above CLT does not depend on the specific form of \(X_t\).
For any bounded local function \(f=f(X_t,\dots,X_{t+\kappa})\), the additive functional under \(\nu_{\mathrm{Parry}}\),
\(
S_n(f):=\sum_{t=1}^{n}f(\sigma^t x)
\),
also satisfies the central limit theorem; its mean density is \(\mathbb E_{\nu_{\mathrm{Parry}}}[f]\),
and the variance density can be obtained from the corresponding Green--Kubo series or the Poisson equation.
Therefore, the \(\varphi\)-closed-form constants describe not only the ``count of \(1\)'s'' but also provide an auditable limiting law template for a class of local observable statistics on the stable layer.
\end{remark}

\paragraph{Tensorization: pullback to general fibers}

\begin{proposition}[Tensorization of hyperplane statistics and local Markov property]\label{prop:pom-fiber-hyperplane-tensorization}
Let \(\mathcal F_x\cong\prod_{i=1}^r\Gamma_{\ell_i}\).
Then:
\begin{enumerate}
  \item \textbf{(Hyperplane count tensorization)}
  For any factor \(j\) and its internal coordinate \(a\in\{1,\dots,\ell_j\}\),
  let \(\Theta_{j,a}\) denote the \(\Theta\)-class corresponding to flipping the \(a\)-th coordinate of the \(j\)-th factor. Then
  \[
  |\Theta_{j,a}|
  =\bigl(F_aF_{\ell_j-a+1}\bigr)\cdot\prod_{i\neq j}F_{\ell_i+2}.
  \]
  \item \textbf{(Bulk probability tensorization)}
  The uniform measure on \(\prod_i V(\Gamma_{\ell_i})\) decomposes strictly as the tensor product of the uniform measures on each factor;
  when \(\ell_i\to\infty\) and the window position remains in the bulk (with distances to both ends tending to infinity),
  the local statistics of each factor converge to the Parry measure by Proposition \ref{prop:pom-unif-xn-local-limit-parry},
  so the product fiber converges locally to an independent tensor product measure.
  \item \textbf{(CLT for additive observables)}
  For component-additive observables (such as the count of ``toggle coordinates equal to \(1\)''), the mean and variance densities are additive over components, and a CLT holds in the bulk limit.
\end{enumerate}
\end{proposition}
\begin{proof}
Part (1) follows from the Cartesian product structure of \(\mathcal F_x\) and the decomposition of \(\Theta\)-classes over products:
for a fixed factor and coordinate direction, the \(\Theta\)-class edge count equals that factor's corresponding \(\Theta\)-class edge count multiplied by the vertex counts of all other factors;
applying Theorem \ref{thm:pom-fibcube-theta-class-size} to each factor yields the result.
Part (2) follows directly from the Cartesian product of vertex sets and the product structure of the uniform measure; the bulk limit is given by Theorem \ref{thm:pom-parry-limit-chain-explicit} and local window convergence.
Part (3) follows from the central limit theorem under independent tensor products and the per-component constants from Theorem \ref{thm:pom-parry-ones-clt}. This completes the proof.
\end{proof}

\begin{corollary}[Invertibility of hyperplane counting: reconstructing component length spectrum from \texorpdfstring{$|V|$}{|V|} and \texorpdfstring{$\{|\Theta|\}$}{|Theta|}]\label{cor:pom-fiber-hyperplane-spectrum-invertibility}
Let \(\mathcal F_x\cong\prod_{i=1}^r\Gamma_{\ell_i}\), and denote \(V_x:=|V(\mathcal F_x)|\) and \(\mathcal M_x:=\{|\Theta|\}\) the multiset of edge counts of all \(\Theta\)-classes of \(\mathcal F_x\).
Then \((V_x,\mathcal M_x)\) uniquely determines the length spectrum multiset
\(
\mathsf{Spec}_{\square}(x)=\{\ell_1,\dots,\ell_r\}
\);
thus on this fiber class, the hyperplane statistics \(\mathcal M_x\) (together with the cardinality \(V_x\)) determine the Cartesian prime factorization and graph isomorphism class of \(\mathcal F_x\) in an implementation-independent sense (Corollary \ref{cor:pom-fiber-reconstruction-spectrum-complete}).
\end{corollary}
\begin{proof}
By Proposition \ref{prop:pom-fiber-hyperplane-tensorization}, for any factor \(j\) and coordinate \(a\in\{1,\dots,\ell_j\}\) we have
\[
\frac{|\Theta_{j,a}|}{V_x}=\frac{F_aF_{\ell_j-a+1}}{F_{\ell_j+2}}.
\]
Therefore the normalized multiset
\(
\mathcal P_x:=\{|\Theta|/V_x:\ |\Theta|\in\mathcal M_x\}
\)
decomposes as the multiset union of the profiles of each factor:
\[
\mathcal P_x=\bigsqcup_{j=1}^r \left\{\frac{F_aF_{\ell_j-a+1}}{F_{\ell_j+2}}:\ 1\le a\le \ell_j\right\}.
\]
For fixed length \(\ell\), the endpoint value
\(
p^{\mathrm{end}}(\ell):=F_{\ell}/F_{\ell+2}
\)
is by Theorem \ref{thm:pom-fibcube-theta-class-size} the unique maximal element of that profile and appears exactly twice (corresponding to \(a=1,\ell\)).
Let \(p_{\max}:=\max\mathcal P_x\); then \(p_{\max}=p^{\mathrm{end}}(\ell)\) for some \(\ell\), and its multiplicity must be \(2m_\ell\), where \(m_\ell:=\#\{j:\ell_j=\ell\}\).
\par
To invert \(p_{\max}\) to obtain \(\ell\), let \(u_n:=F_{n+2}/F_n\) (\(n\ge 1\)).
By the identity
\(
F_{n+3}F_n-F_{n+2}F_{n+1}=(-1)^{n+1}
\),
we have
\[
u_{n+2}-u_n=\frac{(-1)^{n+1}}{F_{n+2}F_n},
\]
so \(u_{2k}\) is strictly decreasing while \(u_{2k+1}\) is strictly increasing, both converging to \(\varphi^2\) from above and below respectively.
Therefore the endpoint ratio \(p^{\mathrm{end}}(n)=1/u_n\) is strictly identifiable as a function of \(n\) (and its parity is determined by comparison with \(\varphi^{-2}\)).
Hence \(\ell\) and \(m_\ell\) are uniquely determined by \(p_{\max}\) and its multiplicity.
\par
After removing \(m_\ell\) copies of the complete length-\(\ell\) profile
\(
\{F_aF_{\ell-a+1}/F_{\ell+2}\}_{a=1}^{\ell}
\)
from \(\mathcal P_x\), the remaining multiset is still the normalized profile union of the remaining factors and can be processed recursively; recursion terminates upon recovering the full length multiset \(\{\ell_i\}\).
Finally, by Corollary \ref{cor:pom-fiber-reconstruction-spectrum-complete}, the graph isomorphism class is determined. \qed
\end{proof}

\begin{remark}[Instantiation at the minimal window \(m=6\)]\label{rem:pom-fiber-hyperplane-m6}
The \(n=6\) row of Table \ref{tab:fibonacci_cube_fvector_small_n} gives \((C(6,0),C(6,1),C(6,2),C(6,3))=(21,38,22,4)\),
where \(C(6,1)=|E(\Gamma_6)|\) simultaneously equals \(\sum_{i=1}^{6}|\Theta_i|\).
By Theorem \ref{thm:pom-fibcube-theta-class-size},
\(
(|\Theta_1|,\dots,|\Theta_6|)=(F_6,F_5F_2,F_3F_4,F_4F_3,F_5F_2,F_6)=(8,5,6,6,5,8)
\),
so the endpoint hyperplanes \(\Theta_1,\Theta_6\) are the unique largest \(\Theta\)-classes.
On the other hand, Corollary \ref{cor:pom-parry-two-point-alternating} gives the alternating correlation decay rate \(\varphi^{-2}\) in the bulk limit,
so that under the ``cubical geometric semantics of time fibers,'' endpoint hyperplanes, the alternating slow mode, and the \(\varphi\)-constant layer align into a single set of purely combinatorial closed-form fingerprints.
\end{remark}
